\documentclass[a4paper,10pt]{report}\input{../0.Préambule/Préambule_cours_prof}\begin{document}

\newpage
\section*{Quotient de nombres relatifs}
\addcontentsline{toc}{section}{Quotient de nombres relatifs}

\vspace{-2mm}
\begin{exer1}
Calcule mentalement :
\begin{enumerate}[font=\color{exer1}]
\begin{tabularx}{\linewidth}{*{5}{>{\arraybackslash}X}}
\item $64 \div (-8)$ & 
\item $42 \div (-6)$ & 
\item $-24 \div (-3)$ & 
\item $81 \div (+9)$ & 
\item $-17 \div (-1)$ \\ 
\item  $-35 \div 7$ & 
\item $(-54) \div (-6)$ & 
\item $25 \div (-5)$ & 
\item $(-4) \div (+4)$ & 
\item $(-29) \div (+1)$\\
\end{tabularx}
\end{enumerate}
\end{exer1}

\begin{exer1}
Sans calculatrice, donne l'écriture décimale de chacun des nombres suivants :

\vspace{-2mm}
\begin{enumerate}[font=\color{exer1}]
\begin{tabularx}{\linewidth}{*{4}{>{\arraybackslash}X}}
\item $\dfrac{-3}{-10}$ & 
\item $-\dfrac{-64}{-8}$ & 
\item $\dfrac{-50}{+100}$ & 
\item $\dfrac{-3}{-2}$\\
\end{tabularx}
\end{enumerate}
\end{exer1}

\begin{exer1}
Complète les quotients sans poser les opérations

\vspace{-2mm}
\begin{enumerate}[font=\color{exer1}]
\begin{tabularx}{\linewidth}{*{4}{>{\arraybackslash}X}}
\item $24 \div \dots ~\dots ~ = -8$ & 
\item $\dots ~ \dots ~ \div 2,5 = -100$ & 
\item $(-24) \div \dots ~\dots ~ = -12$ & 
\item $\dots ~ \dots ~ \div 25 = -5$\\
\item $-18 \div \dots ~\dots ~ = -6$ & 
\item $\dots ~ \dots ~ \div 5 = 100$ & 
\item $25 \div \dots ~\dots ~ = -5$ & 
\item $\dots ~ \dots ~ \div (-1) = 100$\\
\item $-42 \div \dots ~\dots ~ = 6$ & 
\item $\dots ~ \dots ~ \div (-20) = -80$ & 
\item $-16 \div \dots ~\dots ~ = 32$ & 
\item $\dots ~ \dots ~ \div (-7) = 35$\\
\end{tabularx}
\end{enumerate}
\end{exer1}

\begin{exer1}
Compléter le tableau suivant :
\begin{center}
\renewcommand{\arraystretch}{1.7}
\begin{tabularx}{\linewidth}{|*{6}{>{\centering \arraybackslash}X|}}
\hline
\rowcolor{exer1!20}$a$ & $b$ & $c$ & $\dfrac{a}{-b}$ & $(-c) \div b$ & $-\dfrac{c}{-a}$ \\\hline
$-2$& $4$ & $12$ & && \\\hline
$-8$& $-1$ & $-6,4$ &&& \\\hline
$3$& $-1,5$ & $15$ &&& \\\hline
$-1$& $5$ & $-2$ &&& \\\hline
\end{tabularx}
\end{center}		
\end{exer1}

\begin{exer1}

\begin{enumerate}[font=\color{exer1}]
\item Trouve deux nombres relatifs dont le quotient est positif et la somme est négative.
\item Trouve deux nombres relatifs dont le quotient est négatif et la somme est positive.
\item Trouve deux nombres relatifs dont le quotient et la somme sont positifs.
\item Trouve deux nombres relatifs dont le quotient et la somme dont négatifs.
\end{enumerate}
\end{exer1}

\begin{exer1}
Pour chaque quotient, dire si il est positif ou négatif.
\begin{enumerate}[font=\color{exer1}]
\begin{tabularx}{\linewidth}{*{5}{>{\arraybackslash}X}}
\item $\dfrac{12 \times (-2)}{(-4) \times (-8)}$ &
\item $\dfrac{1 \times (-2) \times 3}{4 \times (-7)}$ &
\item $-\dfrac{-2,1}{(-12) \times (-4,2)}$ &
\item $-\dfrac{11 \times (-3)}{(-5) \times (-4)}$ &
%\item $\dfrac{-4 \times 2}{(-5) \times 3}$ &
\item $-\dfrac{11 \times (-3) \times (-2)}{6 \times (-7)}$ \\
\end{tabularx}
\end{enumerate}
\end{exer1}

\begin{exer1}
Calcule les fractions suivantes
\begin{enumerate}[font=\color{exer1}]
\begin{tabularx}{\linewidth}{*{4}{>{\arraybackslash}X}}
\item $A = \dfrac{11 \times (-3)}{(-5) \times (-2)}$ &
\item $B = \dfrac{(-3) \times 2 \times (-5)}{-10 \times 4}$ &
\item $C = \dfrac{7 \times (-2)\times 8}{14 \times 5}$ &
\item $D = \dfrac{(-1) \times (-2) \times (-1)}{5 \times (-4)}$ \\
\end{tabularx}
\end{enumerate}
\end{exer1}

\begin{exer1}
En utilisant la calculatrice, donne une valeur approchée au centième
\begin{enumerate}[font=\color{exer1}]
\begin{tabularx}{\linewidth}{*{6}{>{\arraybackslash}X}}
\item $ (-1) \div 3$ &
\item $ (-5) \div (-11)$ &
\item $ 1,3 \div 0,7$ &
\item $ 2,9 \div (-6)$ &
\item $ 47 \div (-23)$ &
\item $ -9,5 \div 7$ \\
\item $-\dfrac{-53}{16}$ &
\item $\dfrac{-17}{-47}$ &
\item $-\dfrac{-1,7}{-0,7}$ &
\item $\dfrac{11}{-19}$ &
\item $\dfrac{3}{5}$ &
\item $\dfrac{-1}{-7}$ \\
\end{tabularx}
\end{enumerate}
\end{exer1}

\begin{exer2}
On considère ce programme de calcul :

\vspace{-9mm}
\begin{flushright}
\begin{minipage}{8cm}
\begin{edupython}
Choisir un nombre
Augmenter le nombre de -5.
Multiplier le resultat par 4
Soustraire le double du nombre choisi au depart
Diviser le resultat par -2
Ajouter -10
\end{edupython}
\end{minipage}
\end{flushright}

\vspace{-26mm}
\begin{enumerate}[font=\color{exer2}]
\item Exécute le programme avec les nombres :
\begin{enumerate}[font=\color{exer2}]
\begin{tabularx}{\linewidth}{*{6}{>{\arraybackslash}X}}
\item $12$ ?&
\item $-3$ ?&
\item $0$ ? &&&\\
\end{tabularx}
\end{enumerate}
\item Que remarque-t-on ?
\end{enumerate}
\end{exer2}
\end{document}